% !TEX TS-program = lualatex+se+debug
% !TEX encoding = UTF-8 Unicode

\documentclass[12pt, oneside]{extreport}
% -----------------------------------------------------------------------------------
% ------------------------- USTAWIENIA STRONY--------------------------
% -----------------------------------------------------------------------------------
\usepackage{geometry}
\geometry{
	paperwidth=86mm,
	paperheight=160mm,
	left=4mm,
	right=4mm,
	top=4mm,
	bottom=4mm,
	marginparwidth=0mm,
	marginparsep=0mm
	}
\newcommand{\wstawplik}[1]{\input{../include/#1}}%skraca kod
\wstawplik{headerofficium}
% elementy zmienne co do headerofficium:
\graphicspath{{../img/}} % katalog obrazów
%\svgpath{{../include/}}
\grechangestaffsize{16}%rozmiar wyniku GregorioTEX
\grechangestyle{initial}{\greIniZallmanCaps\fontsize{36}{26}\selectfont} % ustawia czcionkę zdefiniowaną w sekcji czcionek o zmniejsza wielkość inicjału do rozmiaru 46pt a baseline do 36pt
\grechangestyle{annotation}{\color{black}\large}% format adnotacji
\gresetlinecolor{black}% Make staff lines red; remove for black:
\grechangedim{spacelinestext}{0.6cm}{scalable}% The space between the lines and the bottom of the text. 	0.60617 cm plus 0 cm minus 0 cm
\grechangedim{spacebeneathtext}{0pt}{scalable} % odstęp poniżej tekstu, domyślnie: 0pt
\grechangedim{abovelinestextraise}{0pt}{scalable} % odstęp poniżej tekstu, domyślnie: 0pt
\gresetlastline{ragged}% justowanie do lewej
% koniec 

% debugowanie:
% %\iffalse %*** %\fi %***
% \tracingall przed problematycznym kodem - bardzo dużo informacji
% lualatex+se+debug - dodatkowe opcje: -file-line-error i -halt-on-error 
% \errorcontextlines
% é á í ó ú ý æ ǽ œ Á, É, Í, Ó, Ú, Ý Æ Ǽ Œ \'{\oe} \'{\OE}
\def\mydebug{1}% 0 - debugowanie, 1 - bez debugowania (wyłącza bloki \if\mydebug1)
\begin{document}
\makeatletter
\grechangedim{interwordspacetext}{\unexpanded{\the\fontdimen2\font plus \the\gluestretch\gre@space@skip@interwordspacenotes minus \the\glueshrink\gre@space@skip@interwordspacenotes}}{fixed}%
\makeatother
\pagestyle{empty}	
\begin{center}
	\begin{latin}
		\subsection*{Si quaeris miracula(Tipica)}
	\end{latin}		
\end{center}
\gresetinitiallines{1}\greannotation{1.}\gabcsnippet{(cb4)Si(d) quæ(fe f!gh)ris(g) mi(h)rá(fvED!e)cu(d)la,(d) (,) mors,(h) er(kj)ror,(k) ca(h)lá(i)mi(i)tas,(h) (,) dæ(k)mon,(j) le(k)pra(h) fú(i)gi(i)unt,(h) (,) æ(f)gri(g) sur(h)gunt(g) sa(fe)ni.(d) *(::) Ce(k)dunt(j) ma(i)re,(h) vín(g)cu(f)la,(f)(,) mem(d)bra,(e) res(f)que(g) pér(hi)di(h)tas(h) (;) pé(d)tunt(e) et(f) ac(g)cí(hi)pi(h)unt(h) (,) jú(i)ve(h)nes(g) et(h) ca(fe)ni.(d)(::) Pé(d)re(e)unt(f) pe(g)rí(hi)cu(h)la,(h) (,) ces(j)sat(iyi) et(j) ne(k)cés(j)si(ixi)tas,(h) (;)nar(f)rent(g) hi(h) qui(i) sén(h)ti(g)unt,(f)(,) di(h)cant(g) Pa(fe)du(d)á(fe)ni.(d) * (::) Ce(k)dunt...(j) (::) Gló(de!fg!hi)ri(h)a(h) Pa(iyji)tri,(j) et(k) Fí(j)li(ixi)o,(h) (,)et(f) Spi(g)rí(hiHG)tu(f)i(f) San(hGFEDe)cto.(d) (::) Ce(k)dunt...(j) (::)}
\end{document}